\paragraph{Curvature at the selected short weight.}
\hypertarget{source-slice-short-analytic}{}

Control on the chosen cut must extend to a neighborhood. The curvature of the short weight determines the local norm comparison used in this extension.

Write \(q\) for the length of the first signature block,
\(p=r-q\), and \(s=fr(r-1)/2\). At an orthonormal flag,
use a lower unipotent coordinate matrix \(I+Z\). The
Gram matrix in each embedding is \((I+Z)^*J(I+Z)\), where
\(J\) is diagonal with \(q\) signs of one kind followed by
\(p\) of the other. The potential of a weight \(v\) is
\[
 \varphi_v=\sum_\tau\sum_{i=0}^{r-1}
               v_{\tau,i}\log\left|
                 \frac{\Delta_{\tau,i+1}}{\Delta_{\tau,i}}\right|,
\]
where \(\Delta_{\tau,k}\) is the leading minor of
order \(k\) and \(\Delta_{\tau,0}=1\).

In the quadratic expansion of this minor's logarithm,
an entry \(Z_{ji}\), \(i<j\), contributes only when
\(i<k\le j\). The first \(k\) rows give a unit lower
triangular factor whose determinant is constant. The
remaining rows contribute the quadratic trace with sign
\(J_j/J_i\). For a tangent connecting positions \(i<j\), the coefficient of the Levi form after summing the weight differences is
\[
  \begin{cases}v_i-v_j&i,j\text{ in the same block},\\v_j-v_i&i,j\text{ in opposite blocks}.\end{cases}
\]
 At another
point take its orthonormal frame by triangular
orthogonalization. The induced change on tangents is
invertible and the logarithm of the change in the weight
character is pluriharmonic. This transports the diagonal
calculation to that point.

In this chart the invariant density is
\(D=\prod_{\tau,k<r}|\Delta_{\tau,k}|^{-2}\).
The full Gram determinant is constant on the unipotent
cell, so

\[
  \mathcal L(\log D)=\mathcal L(\varphi_{k_{\rm vol}})
\]
with \((k_{\rm vol})_{\tau,i}=2i-(r-1)\). Write \(d\)
for the sign weight that is \(-1\) on the first block
and \(+1\) on the second. The radius potential \(R\)
is \(\varphi_d/2\) plus a pluriharmonic function, since
the only nonzero interior weight difference of \(d\)
occurs at \(q\). In particular \(\mathcal L(R)\) has unit
coefficients in the cross-block directions and zero
coefficients within blocks in these coordinates.
\label{cor:slice-short-analytic-1}

\begin{lemma}[Curvature for the short-weight interpolation]
\label{lem:short-weight-curvature}
In the flat cut chart with the given signatures, the short weight
has the positive Levi directions used in the local interpolation
estimate. For the selected cutoff on each fixed zero neighborhood, put \(G=w_F-k_{\rm vol}-s w_L\). With \(i,j\) in opposite blocks in the indicated order, the signature margin is
\[
 G_j-G_i\ge2((w_L)_j-(w_L)_i)\ge19T/5.
\]

\end{lemma}

\begin{proof}
\proofstep{Positivity at the chosen short weight.}
Here is the weight comparison used on both runs.
Put \(G=w_F-k_{\rm vol}-s w_L\). Within either block the
selected weight inequalities give
\(G_i-G_{i+1}\ge2((w_L)_i-(w_L)_{i+1})>0\).
Across the blocks they give
\((w_L)_j-(w_L)_i\ge19T/10\) and
\(G_{r-1}-G_0\ge2((w_L)_{r-1}-(w_L)_0)\).

Summing the adjacent inequalities from the outer
indices to \(i,j\) gives
\[
 G_j-G_i\ge2((w_L)_j-(w_L)_i)\ge19T/5 .
\]
These are the bounds for the actual shifted short weight.
In taking consecutive differences of \(w_F-s w_L\),
the volume weight contributes
\((k_{\rm vol})_{i+1}-(k_{\rm vol})_i=2\). The light weight and
\(2w_L-d\) are strictly positive in the Levi formula.
For \(0\le u\le T\), the weight \(G-u d-w_L/2\)
has nonnegative Levi form. Its cross-block gap is at
least \(3(19T/10)/2-2T>0\), and the within-block
difference of \(d\) is zero.

\proofstep{The weighted curvature comparison.}
The invariant potential \(\psi_0\) is bounded above and
has \(\mathcal L(\psi_0)\ge-(r/2)\mathcal L(\varphi_d)\).
Its cusp coefficient and the selected cutoff are
\[
 a_1=\frac{1}{2\cdot5^4|\operatorname{disc}E|^{1/f}},
 \qquad
 c_1=2(\operatorname{rd}E)^7 T c\,(fq)^2 q .
\]
Here \(q\ge1\) on the indefinite lift and \(c>0\)
is the light first-minimum coefficient. To check the allowance, put \(d_E^{\rm root}=\operatorname{rd}E\), \(\mathcal P=C_2z^{A_2}\), and write \(Q\) for the denominator of the light coefficient.
The numerical inequalities on this run are
\[
 c=(d_E^{\rm root})^7\mathcal P D_F^2/Q,\qquad
 (d_E^{\rm root})^{28}\mathcal P^3D_F^2\le Q,\qquad
 1250\,rfm^3\le\mathcal P^2,\qquad
 2fm\le d_E^{\rm root}.
\]
Thus \(c\mathcal P^2(d_E^{\rm root})^{21}\le1\).

Since \(q\le m\) and \(d_E^{\rm root}\ge1\),
\[
 \frac{r c_1}{a_1}
    =Tc\, (1250\,rfq^3)\,(2f(d_E^{\rm root})^9)
    \le Tc\mathcal P^2(d_E^{\rm root})^{21}\le T .
\]
With \(r\le T\) this gives \(r(1+c_1/a_1)\le2T\).
Put \(\psi=(1+c_1/a_1)\psi_0\). The preceding
comparison with \(u=r(1+c_1/a_1)/2\) gives
\[
 \mathcal L(\varphi_F-\log D-s\varphi_L+\psi)
                    \ge \tfrac12\mathcal L(\varphi_L).
\]
The cross-block estimate also gives the independent
lower bound \(m_{\rm int}\mathcal L(R)\) on this expression,
with \(m_{\rm int}=19T/5-2T\).

For division the potential
omits the \(\psi\) term and the subtraction of
\(s\varphi_L\). The margin there is
\(m_{\rm div}=19T/5+s(19T/10)\).
Set \(\beta=5fr^2\) and \(u_0=f\min(p,q)\). Then
\[
 4fr^2<\beta<m_{\rm int},
 \qquad \beta^2 u_0+\beta < m_{\rm div}.
\]
Indeed \(f\ge1\), \(r\ge8\), \(2u_0\le fr\), and
\(T\ge11s^2\). These give
\(\beta\le20s^2/49\) and
\(\beta^2 u_0\le20s^3\) by \(5r^2\le(r-1)^3\).
The additional \(19T/5\) and \(19sT/10\) cover the
two terms in the division margin. For a definite
signature \(R=0\) and \(\psi_0=0\); the within-block
calculation supplies the nonnegative Levi form and the
same inequalities with arbitrary margins in the zero
directions of \(\mathcal L(R)\).
\label{cor:slice-short-analytic-2}

\end{proof}

\paragraph{The Poincar\'e estimate on coordinate balls.}

The curvature estimate is local. To obtain the norm bound after Poincar\'e summation one must compare the coordinate balls about the translated frames.

Let \(x_*\) be the orthonormal frame of the seed cell.
At an orthonormal frame \(x\) write the relative
isometry \(x_*^{-1}x\) at each embedding in signature
order \(p,q\), with bottom-right block \(D_0\).
Then \(R(x)=2\sum_\tau\log |\det D_{0,\tau}|\).
For a general admissible frame, subtract the logarithm
of its Gram minor of order \(q\) in each embedding.
This formula is invariant under the upper triangular
change of frame.

\begin{lemma}[Local integrals and the Poincaré estimate]
\label{lem:short-poincare-summation}
For the coordinate balls and weights below, the estimate for the
Poincaré sum follows from the local integral of the lifted section
and the bound on the overlap of the translated balls. In particular, for a seed of finite weighted integral \(I=\int|P|^2W_u e^{\beta R}D\), its integral \(I_\gamma\) on a translated ball, before the factor \(e^{\beta R}D\), satisfies
\[
 I_\gamma\le C I e^{-\beta R(\gamma x)}.
\]
The constants may depend on the fixed cut chart.
\end{lemma}

\begin{proof}
\proofstep{Bounding the overlap of translated balls.}
Use discreteness of the group in the product of the
archimedean matrix groups. Two distinct relative
isometries \(g,h\) from \(\Gamma x\) have an entry of
\(hg^{-1}-I\) bounded below in absolute value by a
positive constant depending on the group and \(x_*\).
This is independent of \(x\) because the factor \(x\)
cancels in the quotient. If all entries of both
isometries have absolute value at most \(A\ge1\),
the entries of \(g^{-1}=Jg^*J\) are at most \(A\).
Since 
\[
  hg^{-1}-I=(h-g)g^{-1},
\]
 at least one entry of
\(h-g\) is separated by a constant times \(A^{-1}\).
Partitioning the box of real and imaginary matrix
entries into cubes of that sidelength bounds its
number of isometries by \(C A^{4fr^2}\).

For a single signature matrix \(g=\left(\begin{smallmatrix}
A_0&B_0\\ C_0&D_0\end{smallmatrix}\right)\) write \(s_i\) for the singular values of \(B_0\). The identity
\(D_0^*D_0=I+B_0^*B_0\) gives
\[
 |\det D_0|^2=\prod_i(1+s_i^2)\ge1+\sum_i s_i^2
 =1+\|B_0\|_{\mathrm{HS}}^2.
\]
The other block identities bound the square of
each entry of \(g\) by \((r+1)|\det D_0|^2\).
Thus a dyadic shell of maximal entry \(A\) contributes
\(O(A^{4fr^2-2v})\) to the sum of \(e^{-vR}\).
Consequently
\[
          \sup_{x\ {\rm orthonormal}}\
               \sum_{\gamma\in\Gamma}e^{-vR(\gamma x)}<\infty
                 \quad(v>2fr^2).
\]

The same bound controls local integrals. On a small
fixed unipotent coordinate ball about an orthonormal
centre \(y\), the Gram minors are close to their identity
values, uniformly in \(y\). In signature order, right
multiplication of the relative isometry by the nearby
coordinate frame \(N\) changes its bottom-right
determinant by the factor
\(\det(N_{22}+D_0^{-1}C_0N_{12})\).
The relation \(D_0D_0^*=I+C_0C_0^*\) implies that \(D_0\) is
invertible and gives the operator inequality
\[
 (D_0^{-1}C_0)(D_0^{-1}C_0)^*
 =I-D_0^{-1}(D_0^{-1})^*\le I.
\]
Hence \(\|D_0^{-1}C_0\|_{\rm op}\le1\). On a ball fixed in
these coordinates the determinant of the factor is at
least \(1/2\), independently of the relative isometry.
The Gram-minor formula gives \(R(yN)\ge R(y)-3f\log2\).
The density \(D\) is bounded below there.

\proofstep{Weights in the local integral estimate.}
Let \(W_u\) be a squared-reading weight transforming
inversely to the squared weight character and invariant
under \(\Gamma\). Suppose the seed has finite integral
\(I=\int |P|^2 W_u e^{\beta R}D\) on the entire flag
cell. Write \(I_\gamma\) for the integral without
\(e^{\beta R}D\) in the coordinates of the translated
ball. The invariant measure and the uniform bounds
give 
\[
  I_\gamma\le C I e^{-\beta R(\gamma x)}.
\]
Here the change to the seed cell preserves the
integral even when the small ball crosses its chart
divisor, which has measure zero. Summing square roots
is allowed since \(\beta/2>2fr^2\).
The holomorphic submean inequality gives the bound on
the normalized Poincar\'e sum by a constant times
\[
       \sup_{z\ \text{in the ball}}
              \left(e^{-\varphi_v(z)}/W_u(z)\right)^{1/2}.
\]

On a smaller neighborhood of any centre the same
submean inequality supplies a common summable
majorant for all evaluation points, provided
\(W_u\ge c_*e^{-\varphi_v}\) for some \(c_*>0\).
Triangular orthogonalization at that centre changes
this majorant by its nonzero weight character. This
proves locally normal convergence on the frame domain.
All constants in the ball bound can depend on the
lift, weight, cut and seed. They are uniform in the
orthonormal centre and its orbit, which is the
uniformity needed for the Fourier test.
\label{cor:slice-short-analytic-3}

\end{proof}

\paragraph{Extracting Fourier coefficients along the cusp flow.}
On an indefinite component the potential is the
logarithm of the sum of normalized squared theta
values over ordered \(r\)-tuples of points of order
dividing 25 whose multiples by 25 generate the
period lattice modulo 5 over \(\mathcal O_E\).
Each factor is divided by the theta squared
\(L^2\) norm. The theta estimate used here assumes an
integral CM element \(\omega\) with
\(\omega+\bar\omega=1\), the even trace comparison for
multiplication by \(\omega-\bar\omega\), an isotropic
lattice vector, \(\operatorname{rd}E\ge\max(f^2,16)\),
and a rank parameter \(h\) with \(h^4\le f\) and
\(\min(p,q)\le h\). In the trace comparison the difference
of the two trace pairings on any two integral period
vectors is in \(2\mathbb Z\). The chosen host supplies
\(\omega\) and this comparison.

On lift \(j\) the
signature has 
\[
  q=m-j>0,
\]
 so one can use \(h=m\).
The invariant smooth potential satisfies the lower Levi
bound used above and, for \(0<\varepsilon<1/2\),
\[
 \psi_0(x)\le C_\varepsilon
    +A_\varepsilon\log(1+\lambda_1(x)^{-1})
    -8\pi(1-2\varepsilon)a_1\min_\tau t_\tau(x,\xi)
\]
whenever all \(t_\tau(x,\xi)\ge1\) for an isotropic
lattice vector \(\xi\). The bound with just the constant
and logarithm also holds on the entire domain.
Here \(\lambda_1^2\) is the minimum of the sum of the
positive majorants on nonzero lattice vectors, and
\(t_\tau(x,\xi)\) is the reciprocal of the corresponding
individual weighted majorant. In the thin region a
vector with sum of the majorants at most
\(2\lambda_1^2\) is isotropic by integrality and all its
times are at least \((2\lambda_1^2)^{-1}\). Using it in
the displayed inequality bounds \(\psi_0\) above even
as \(\lambda_1\) tends to zero.
\label{cor:slice-short-analytic-4}

\begin{proposition}[Fourier vanishing and the cusp decay estimate]
\label{prop:slice-cusp-gap-decay}
At the prescribed short weight the vanishing test on each complete rank-one period forces its low Fourier--Jacobi coefficients \(a_\nu\) to satisfy
\[
 a_\nu=0\qquad\left(\sum_\tau\nu_\tau\le c_1\right).
\]
 Here \(\tau\) ranges over the chosen CM half, so
\(\sum_\tau\nu_\tau=\tfrac12\operatorname{Tr}_{E/\mathbb Q}(\nu)\).
The resulting bound on a section with those terms zero
is the uniform first-minimum decay used in the division norm.
The cross-block gap gives the stated limiting constant-coefficient
test also for \(w_F+k w_L\) with \(k\ge0\).
\end{proposition}

\begin{proof}
\proofstep{The local estimate for coefficient tests.}
For the coefficient test, a small coordinate ball
about any orthonormal flag changes each positive
majorant by factors in \([1-\delta,1+\delta]\), for
any fixed \(\delta>0\) after decreasing the radius.
Indeed the majorant matrices in that ball are given by
orthogonalizing the Gram matrix of the unipotent
coordinate frame. In orthonormal coordinates these
matrices vary continuously from the fixed positive
form at the origin. The comparison holds simultaneously
for all vectors, so it gives the same comparison for
the squared minima and a reciprocal comparison for
each time. Put \(a=(1+c_1/a_1)a_1=a_1+c_1\).

For every \(a'<a\) with \(a'>0\), choose the radius
and \(\varepsilon\) so that
\((1-2\varepsilon)a/(1+\delta)>a'\).
The normalized ball bound for the Poincar\'e sum then
gives
\[
 \|P_{\rm aut}(x)\|_v
     \le C(1+\lambda_1(x)^{-1})^A
               e^{-4\pi a'\min_\tau t_\tau(x,\xi)}
\]
when all times at the centre are at least 2. Triangular
normalization carries it to general frames with the
same majorants. The exponent \(4\pi\) comes from the
square root of the metric bound.

\proofstep{Decay along the integral cusp flow.}
On the holomorphic flow of \(\xi\), translate each
complex parameter to \(u_\tau+iY\), with the real
parameters \(u_\tau\) in a fixed period parallelepiped.
Each time becomes \(t_\tau(x,\xi)+Y\). This follows
from the rank-one isotropic shear. The Gram matrix at
time \(Y\) is a rank-one Hermitian update of the
original Gram matrix, since the quadratic term of the
shear is a multiple of the norm of \(\xi\) and vanishes.
Each leading minor divided by its value at zero is
\(1+Y b_{\tau,k}\) with \(b_{\tau,k}\ge0\).
Writing \(\Delta_{\tau,k}(Y)\) for the moving leading minor, the rank-one calculation gives
\[
  \frac{\Delta_{\tau,k}(Y)}{\Delta_{\tau,k}(0)}=1+Yb_{\tau,k}\ge1\qquad(Y\ge0).
\]

The signs of the minors persist for all \(Y\ge0\).
This gives polynomial upper and inverse bounds for
the diagonal factors of the weight metric. The entries
of the frame and of its inverse grow at most polynomially
in \(Y\) for bounded \(u\), since the shear is linear
in the flow parameter and its rank-one term squares
to zero. Orthogonalization using these minors shows
that the inverse of the positive majorant matrix grows
at most polynomially as well. Hence
\(1+\lambda_1^{-1}\) and the change from the normalized
norm to the holomorphic reading both have polynomial
bounds in \(1+Y\) on this slice.
In particular along this slice there are constants \(C,N\), depending on the initial data, for which
\[
  1+\lambda_1(Y)^{-1}\le C(1+Y)^N .
\]
 These bounds may depend
on the initial flag and on \(\xi\).

The holomorphic reading has character frequency
\(2\nu\) on the real period lattice. By definition the
coefficient is zero when \(2\nu\) is outside the dual
period lattice. For an actual character, shifting the torus
integral to height \(Y\), by holomorphy and periodicity,
gives
\[
 |F_{\xi,\nu}(x)|e^{-4\pi Y\sum_\tau\nu_\tau}
     \le \sup_u |P_{\rm aut}(\operatorname{flow}_\xi(u+iY)x)|.
\]
If \(\sum_\tau\nu_\tau<a\), take \(a'\) strictly between
\(\max(0,\sum_\tau\nu_\tau)\) and \(a\). The right hand
side is bounded by a polynomial times
\(e^{-4\pi a'Y}\). Letting \(Y\) grow gives the required
vanishing. The strict inequality at \(a=a_1+c_1\) kills
the mode of total frequency exactly \(c_1\).
\label{cor:slice-short-analytic-5}

\paragraph{Uniform decay of a section with zero low modes.}
\label{par:slice-cusp-decay}

The extracted low Fourier coefficients describe the part that must be removed. Once those coefficients vanish, the cusp flow bounds the remainder uniformly.

Let \(P\) be a bounded normalized section whose coefficients
of total frequency at most \(c_1\) vanish on the rank-one
periods. The moderate-growth coefficient estimate, using
one quarter of its exponential rate as slack, gives
\[
 \|P_{\xi,\nu}(x)\|_\kappa
 \le C \exp\!\left(-3\pi\sum_\tau\nu_\tau
                                      t_\tau(x,\xi)\right)
\]
when the majorants of the integral isotropic vector \(\xi\)
are at most one. The constant is independent of the vector,
the frequency and the flag \(x\) on the specified quotient.
The support has \(\nu_\tau\ge0\), and the nonzero terms are
admissible frequencies. The frequency-sum bound on the
period lattices gives a uniform finite majorant for the sum
of \(e^{-(3\pi/2)\sum_\tau\nu_\tau}\) over their nonnegative
frequencies. These coefficient and summability bounds apply
to the bounded weight-\(\kappa\) sections here.

\proofstep{Using the first minimum to choose the flow.}
For \(\lambda_1(x)^2\) sufficiently small, choose an integral
vector with sum of its weighted majorants

\[
  S<3\lambda_1(x)^2/2
\]
 and \(S<1\).
At this threshold the vector is isotropic, by integrality
of its Hermitian norm and the bounds under all embeddings.
Its times are at least \(1/S\). Each surviving \(\nu\) has
\(\sum_\tau\nu_\tau>c_1\). Half of the exponential rate
therefore bounds the term by \(C e^{-\pi c_1/\lambda_1^2}\),
and the other half leaves the summable frequency majorant
above. The period expansion and the triangle inequality
give \[
 \|P(x)\|_\kappa\ll e^{-\pi c_1/\lambda_1(x)^2}.
\]
On the rest of the quotient the positive lower threshold
for \(\lambda_1(x)\) and boundedness give the same estimate
with a larger constant. The squared norm bound is uniform
in \(x\) at the given lift. Its constant can depend on the
lift and section since it is used for integrability.
\label{cor:slice-short-analytic-6}

\proofstep{The norm needed by division.}
For the division coefficients only the zero mode must
vanish. A bounded holomorphic section of a weight \(v\)
with nonincreasing weights inside each block and

\[
  v_{\tau,r-1}-v_{\tau,0}\ge1
\]
 at one embedding \(\tau\)
is already cuspidal. Shift the imaginary flow at this
embedding alone. Let \(\rho_k\) be the ratio of the
leading Gram minor at the new time to its original
value. The isotropic rank-one update has
\(\rho_0=\rho_r=1\) and \(1\le \rho_k\le\rho_q\):
the signed squared lengths in the isotropic vector's
orthogonal-frame coordinates accumulate to an extreme
at the cut between signs. Also 
\[
  \rho_q=1+Y/t_\tau.
\]
Summation by parts yields
\[
 \sum_{i=0}^{r-1} v_{\tau,i}
       (\log\rho_{i+1}-\log\rho_i)
     \le (v_{\tau,0}-v_{\tau,r-1})\log\rho_q .
\]
The coefficients of \(\log\rho_k\) for \(k\ne q\) are
nonnegative and the term at \(q\) keeps its exact value.

The metric conversion now bounds the holomorphic
reading on the translated period torus by a constant
times \((1+Y/t_\tau)^{-1/2}\).
The other embeddings retain their metric factors.
The torus integral for the zero coefficient tends
to zero. These weight conditions hold for \(w_F+k w_L\)
with \(k\ge0\). In a block the integral light gap is
at least one, so the two light gaps in \(G\) absorb
the successive change of two in \(k_{\rm vol}\). The
outer cross-block gap is at least one by the bounds
on \(G\), \(w_L\) and the increasing \(k_{\rm vol}\).
\label{cor:slice-short-analytic-7}
\end{proof}
